

In this section, we focus on the problem of sampling from a Gaussian tilt of the form
\begin{align}\label{eq:RGO}
    \rgo(x)\propto \exp\prn[\Big]{-f(x)-\frac{\nrm{x-\xz}^2}{2\eta}}\,,
\end{align}
where we assume access to first-order queries for $f$. For our eventual application to diffusion models, we will choose $f=\scedit{-}\log p_t$ in view of \cref{eq:backward_kernel} and approximate $\scedit{-}\nabla f\approx \scf_t$. 

\subsection{First-order rejection sampling (FORS)}\label{sec:motivation}




We motivate our approach by considering the simple problem of sampling from a density $p\propto e^{-f}$, where $f : [0,1]\to\R$, $f(0) = 0$, and $-1 \le f' \le 1$.
Our goal is to develop a \emph{high-accuracy sampler}---a sampler whose sampled distribution has $\delta$ error in total variation distance to $p$---in $\polylog(1/\delta)$ steps, using only queries to $f'$.

Consider performing rejection sampling with the base measure $\unif([0,1])$. To do so, we must generate randomness $b\sim \Ber(ce^{-f(x)})$ for any given $x\in[0,1]$. However, we only have access to $f'$. The identity $f(x)=\int_{0}^x f'(y)\,dy$, written as
\begin{align}
\label{eq:1d_representation}
    f(x)=\En_{y\sim \unif([0,x])}[xf'(y)]\,,
\end{align}
suggests an unbiased estimate of $f(x)$. Is it possible to sample from $\Ber(ce^{-f(x)})$ if we have access to an unbiased estimate of $f(x)$?
A more general version of this idea is known as the ``Bernoulli factory'' problem \citep{keane1994bernoulli,nacu2005fast}. 
It can be stated as the following abstract task:
\begin{center}
    \textbf{Task:} Given i.i.d.\ random variables $W_1,W_2,W_3,\dotsc$ in $[-1,1]$, generate a sample $b\sim \Ber(ce^{\En W_1})$.
\end{center}
To solve this, we write the Taylor series as
\begin{align}\notag
    e^{\En W_1}=e^{-1}\cdot e^{\En[1+W_1]}
    =\sum_{j\geq 0} \frac{e^{-1}}{j!}\, \bigl(\En[1+W_1]\bigr)^j\,.
\end{align}
Suppose that $J\sim \Poi(2)$ is independent of the i.i.d. sequence $W_1,W_2,W_3,\dotsc$. Then we notice that
\begin{align}\notag
    e^{\En W_1}=e\En\brk[\Big]{\prod_{j=1}^J \bigl(\frac{1+W_j}{2}\bigr)}\,,
\end{align}
and simply set $b\sim \Ber\prn[\big]{\prod_{j=1}^J \bigl(\frac{1+W_j}{2}\bigr)}$.
Indeed, $\Pr(b=1) = \E\prod_{j=1}^J \bigl(\frac{1+W_j}{2}\bigr) = e^{-1+\En W_1}$.

In summary, we can generate a sample $b\sim \Ber\prn{ce^{-f(x)}}$ \emph{without} having to compute or accurately approximate $f(x)$ (this requires the integration step $f(x)=\int_{0}^x f'(y)\,dy$ and can be expensive). Instead, it is sufficient to have access to (a random number of) unbiased estimates of $f(x)$. The latter can be achieved with derivative information, in view of \eqref{eq:1d_representation}.

We now generalize this setup via the following meta-algorithm, called \emph{first-order rejection sampling} (\emph{FORS}). Given a proposal distribution $q$ and a tilt function $w$, the goal of \cref{alg:fors} is to produce a sample from $\widehat p(x)\propto q(x)\, e^{w(x)}$ without having access to the \emph{value} $w(x)$. Instead, for each $x\in\RR^d$, we can generate i.i.d.\ samples $W_1,W_2, W_3\dotsc$ such that $\En[W_1\mid x]=w(x)$. 
Let $\mathcal W_x$ denote the conditional distribution of $W_1$ given $x$.

\begin{algorithm}
\caption{First-order rejection sampling (FORS)}\label{alg:fors}
\begin{algorithmic}
\STATE \textbf{Input:} Parameter $B > 0$, proposal distribution $q$ over $\R^d$, estimator distributions $(\cW_x)_{x\in\R^d}$ supported on $[-B,B]$
\FOR{$i=1,2,3,\dotsc$}
\STATE Sample $x\sim q$.
\STATE Sample $J\sim \Poi(2\B)$.
\STATE Sample i.i.d. $W_1,\dotsc,W_J \sim \mathcal W_x$.
\STATE Output $x$ with probability $\prod_{j=1}^J \frac{B+W_j}{2B}$.
\ENDFOR
\end{algorithmic}
\end{algorithm}


\begin{theorem}[FORS guarantee]\label{thm:fors}
    \cref{alg:fors} outputs a random point with density $\widehat p(x)\propto q(x)\, e^{\En[W_1\mid x]}$.
    The number of sampled $W_j$'s is bounded, with probability at least $1-\delta$, by $3Be^{2B}\log(2/\delta)$.

    Moreover, if~\cref{alg:fors} is called $T$ times, then with probability at least $1-\delta$, the total number of sampled $W_j$'s is $O(Be^{2B}\,(T+\log(1/\delta)))$.
\end{theorem}

\fcedit{We remark that variants of this idea have been applied to \emph{exactly} simulate SDEs~\citep[e.g.,][]{Wag1988SDE,beskos2005exact,beskos2006retrospective, Pap11Diffusion}.}


\subsection{Sampling from Gaussian tilts with an exact oracle}

Now, we return to the problem of sampling from a general Gaussian tilt:
\begin{align*}
    \rgo(x)\propto \exp\prn[\Big]{-f(x)-\frac{\nrm{x-\xz}^2}{2\eta}}\,.
\end{align*}
\fcedit{From \cref{sec:motivation}, it suffices to construct a proposal distribution $q$ and a tilt function $w$ such that (a) $\nu(x)\propto q(x)\cdot e^{w(x)}$, and (b) a bounded estimator for $w(x)$ can be constructed from $\nabla f$.

The condition (a) is equivalent to
\begin{align*}
    w(x)=-f(x)-\frac{\nrm{x-\xz}^2}{2\eta}- \log q(x)+\const.
\end{align*}
To ensure both (a) and (b), the natural idea is to choose $q$ as a Gaussian approximation to $\nu$ obtained via a first-order expansion of $f$.
Concretely, by Taylor's expansion, for a fixed  $\xp\in\RR^d$, we have
\begin{align*}
    f(x)\approx f(\xp)+\tri{x-\xp, \nabla f(\xp)}\,,
\end{align*}
and hence it is natural to choose $q=\normal{\xz-\eta \nabla f(\xp),\eta\Id}$ so that
\begin{align*}
    q(x) \propto \exp\Bigl(-\tri{x-\xp, \nabla f(\xp)} - \frac{1}{2\eta}\norm{x-x_0}^2\Bigr)\,.
\end{align*}
}
Then, we can express
\begin{align*}
    \icml{&~}\log \rgo(x)-\log q(x)-\const \icml{\\}
    =&~\tri{x-\xp, \nabla f(\xp)}-f(x) + f(x_+) \\
    =&~\int_{0}^1 \tri{x-\xp, \nabla f(\xp)-\nabla f(\lr x+(1-\lr)\xp)}\,d\lr\,.
\end{align*}
In particular, we can set
\begin{align*}
    W_{r,x}\ldef \tri{x-\xp, \nabla f(\xp)-\nabla f(\lr x+(1-\lr)\xp)}\,,
\end{align*}
so that
\begin{align*}
    \rgo(x)\propto q(x)\exp\prn*{ \En_{r} W_{r,x}}\,,
\end{align*}
where the expectation $\En_r[\cdot]$ is taken over $r\sim \unif([0,1])$. 

Further, assuming that $\nabla f$ is $\beta$-Lipschitz, we can bound $\abs{W_{r,x}}\leq \beta\,\nrm{x-\xp}^2$. Therefore, as long as $\xp\approx \xz-\eta\nabla f(\xp)$ and $\beta d\eta\ll 1$, we can guarantee that $\sup_{r}{\abs{W_{r,x}}}\leq 1$ with high probability over $x\sim q$. This implies the following distribution is a good approximation of $\rgo$:
\begin{align*}
    \wh{\rgo}(x)\propto q(x)\exp\prn{\En_{r} \hW_{r,x}}\,,
\end{align*}
where $\hW_{r,x}=\pclip{W_{r,x}}$ and $B=\Theta(1)$ is a tuneable parameter.

\paragraph{General path integral}
We can generalize the above pipeline with any \emph{path integral} and demonstrate (a) that in the gradient Lipschitz case, we in fact only need to ensure $\beta\sqrt{d} \eta\ll 1$, and (b) the Lipschitz continuity of $\nabla f$ can be weakened to H\"older continuity for any exponent $s \in [0,1]$.

Fix a distribution $P$ (to be determined later) over $\RR^d$ and a \emph{path function} $\gam{\lr}\ldef \gamma(x;z,r)$ such that $\gam{1}=x$ and $\gam{0}=\wb{\gamma}(z)$ is independent of $x$. Then we can express any smooth function $h:\RR^d\to \RR$ as
\begin{align}\label{eq:path-integral}
\begin{aligned}
    \icml{&~}h(x)-\En_{z\sim P}[h(\wb{\gamma}(z))] \icml{\\}
    =&~\int_{0}^1 \En_{z\sim P}\tri{\gamp{\lr}, \nabla h(\gam{\lr})}\,d\lr \\
    =&~\En_{z\sim P,\,r\sim \unif([0,1])}\tri{\gamp{\lr}, \nabla h(\gam{\lr})}\,,
\end{aligned}
\end{align}
where $\gamp{\lr}=\frac{d}{d\lr}\gam{\lr}$ is the path derivative with respect to $\lr$. 
This inspires us to consider
\begin{align}\label{eq:Wrzx}
    W_{r,z,x}\ldef \tri{\gamp{\lr}, \nabla f(\xp)- \nabla f(\gam{\lr})}\,,
\end{align}
so that
\begin{align*}
    \rgo(x)\propto q(x)\exp\prn*{ \En_{r,z} W_{r,z,x}}\,,
\end{align*}
where the expectation $\En_{r,z}[\cdot]$ is taken over $r\sim \unif([0,1])$ and $z\sim P$. 
Again, we will truncate $W_{r,z,x}$ to ensure that the estimator lies in $[-B,B]$.

For better dimension dependence, we consider the following path function ($\xhat=\xz-\eta\nabla f(x_+)$): 
\begin{align}\label{eq:def-path}
\begin{aligned}
    &~\gam{\lr}=\alr x + (1-\alr) \xhat + \blr z\,, \\
    &~\alr=\sin(\pi\lr/2)\,, ~~\blr=\cos(\pi\lr/2)\,,
\end{aligned}
\end{align}
so that $\gamp{\lr}=\alrp (x -\xhat) + \blrp z$. 

\begin{assumption}\label{ass:holder}
    There exists $s \in [0,1]$ and $\beta_s \ge 0$ such that $\norm{\nabla f(x) -\nabla f(y)} \le \beta_s\norm{x-y}^s$ for all $x,y\in\R^d$.
\end{assumption}

\begin{theorem}\label{thm:gaussian_tilt}
    Suppose that~\cref{ass:holder} holds, $B=\Theta(1)$, and $\nrm{\xz-\eta\nabla f(x_+) - x_+} \leq (d\eta)^{1/2}$.
    
    Consider instantiating~\cref{alg:fors} with the choices $q = \normal{\xz-\eta \nabla f(\xp),\eta\Id}$, and $\mathcal W_x$ the law of $\pclip{W_{r,z,x}}$, where $W_{r,z,x}$ is defined in~\eqref{eq:Wrzx}-\eqref{eq:def-path} and $z\sim \normal{0,\eta\Id}$, $r\sim \unif([0,1])$. 
    Then, the law $\wh\nu$ of~\cref{alg:fors} satisfies $\Dchis{\nu}{\wh\nu} \le \delta^2$, provided that
    \begin{align*}
        \eta^{-1} \gg \prn[\Big]{ \beta_s^2d^s\log(1/\delta)+\frac{s\beta_s^2}{d^{1-s}}\log^2(1/\delta) }^{1/(1+s)}.
    \end{align*}
\end{theorem}

Note that $x_0 - \eta\nabla f(x_+) - x_+ = 0$ when $x_+ = \text{prox}_{\eta f}(x_0)$, that is, the requirement of~\cref{thm:gaussian_tilt} is that we can approximately take a proximal step on $f$ from $x_0$.

\cref{thm:gaussian_tilt} interpolates between the Lipschitz case $s=0$, which requires $\eta^{-1} \gg \beta_0^2 \log(1/\delta)$, and the smooth case $s=1$, which requires $\eta^{-1} \gg \beta_1\, d^{1/2} \log(1/\delta)$.
This recovers the result of~\citet{FanYuaChe23ImprovedProx}, except that we only use first-order queries; we spell out the implications for log-concave sampling in \cref{sec:log_concave}.
